% Options for packages loaded elsewhere
\PassOptionsToPackage{unicode}{hyperref}
\PassOptionsToPackage{hyphens}{url}
\PassOptionsToPackage{dvipsnames,svgnames,x11names}{xcolor}
%
\documentclass[
]{article}
\usepackage{amsmath,amssymb}
\usepackage{iftex}
\ifPDFTeX
  \usepackage[T1]{fontenc}
  \usepackage[utf8]{inputenc}
  \usepackage{textcomp} % provide euro and other symbols
\else % if luatex or xetex
  \usepackage{unicode-math} % this also loads fontspec
  \defaultfontfeatures{Scale=MatchLowercase}
  \defaultfontfeatures[\rmfamily]{Ligatures=TeX,Scale=1}
\fi
\usepackage{lmodern}
\ifPDFTeX\else
  % xetex/luatex font selection
\fi
% Use upquote if available, for straight quotes in verbatim environments
\IfFileExists{upquote.sty}{\usepackage{upquote}}{}
\IfFileExists{microtype.sty}{% use microtype if available
  \usepackage[]{microtype}
  \UseMicrotypeSet[protrusion]{basicmath} % disable protrusion for tt fonts
}{}
\makeatletter
\@ifundefined{KOMAClassName}{% if non-KOMA class
  \IfFileExists{parskip.sty}{%
    \usepackage{parskip}
  }{% else
    \setlength{\parindent}{0pt}
    \setlength{\parskip}{6pt plus 2pt minus 1pt}}
}{% if KOMA class
  \KOMAoptions{parskip=half}}
\makeatother
\usepackage{xcolor}
\usepackage[margin=1in]{geometry}
\setlength{\emergencystretch}{3em} % prevent overfull lines
\providecommand{\tightlist}{%
  \setlength{\itemsep}{0pt}\setlength{\parskip}{0pt}}
\setcounter{secnumdepth}{-\maxdimen} % remove section numbering
\usepackage[most]{tcolorbox}
\usepackage{makeidx}
\makeindex
\ifLuaTeX
  \usepackage{selnolig}  % disable illegal ligatures
\fi
\IfFileExists{bookmark.sty}{\usepackage{bookmark}}{\usepackage{hyperref}}
\IfFileExists{xurl.sty}{\usepackage{xurl}}{} % add URL line breaks if available
\urlstyle{same}
\hypersetup{
  colorlinks=true,
  linkcolor={Maroon},
  filecolor={Maroon},
  citecolor={Blue},
  urlcolor={Blue},
  pdfcreator={LaTeX via pandoc}}

\usepackage{etoolbox}
\makeatletter
\providecommand{\subtitle}[1]{% add subtitle to \maketitle
  \apptocmd{\@title}{\par {\large #1 \par}}{}{}
}
\makeatother
\subtitle{April 03, 2026}
\author{}
\date{}

\begin{document}

\hypertarget{c-for-c-programmers}{%
\section{Gorgo C for C++ Programmers}\label{c-for-c-programmers}}

\newpage

\hypertarget{how-to-use-this-booklet}{%
\section{0. How to use this booklet}\label{how-to-use-this-booklet}}

This is a short booklet to help a C++ programmer get up to speed with C.
As a side benefit, you will probably find that you understand C++ better
at the end of it.

Each chapter is meant to help you understand a topic, but you will still
want to reference API descriptions for more specifics of the parameters
and operating conditions. Hopefully, you'll have the context you need to
understand API documentation.

\hypertarget{tips}{%
\subsection{Tips}\label{tips}}

\textbf{Tips} call out details that you need to pay special attention
to. \textbf{Traps} warn you of common mistakes made. \textbf{Wut} calls
out a detail that is counter-intuitive, so make sure you pay attention.

\hypertarget{try-it}{%
\subsection{Try It}\label{try-it}}

As the intro to the most amazing programming language book ever written
starts out:

\begin{quote}
The only way to learn a new programming language is by writing programs
in it.
\end{quote}

You need to write some code. Make sure you try writing some programs
from scratch. At the end of most sections is a starter program that you
can type in and modify to play with. Don't use it as an excuse to avoid
writing some of your own starter programs. It's the only way to master a
language.

\hypertarget{exercises}{%
\subsection{Exercises}\label{exercises}}

Don't skip the exercises at the end of the chapters. You can get the
answer key, but don't look at the answer key before you work out the
answer yourself. If you look at the answer key first, the concepts will
not sink in.

\end{document}
